% Os três barramentos e o que corre em cada um.
%
% Os três são desenhados dentro de uma mesma faixa porque, no equipamento, eles
% são o mesmo caminho físico: um feixe de vias paralelas ao qual todos os blocos
% se ligam. Desenhá-los como três linhas separadas, cada uma com sua descida,
% obrigaria as descidas a cruzar linhas às quais não se ligam.
%
% O sentido das setas é a informação principal da figura: o de endereços parte
% sempre do processador, e os outros dois têm tráfego nos dois sentidos.
\documentclass[border=5pt]{standalone}
\usepackage{figuras}
\begin{document}
\begin{tikzpicture}

  % --- Quem está ligado ao barramento ------------------------------------
  \node[bloco proc, minimum width=34mm, minimum height=13mm] (cpu) at (0, 0)
    {Processador};
  \node[bloco mem, minimum width=34mm, minimum height=13mm] (mem) at (4.4, 0)
    {Memória\\principal};
  \node[bloco es, minimum width=34mm, minimum height=13mm] (cd) at (8.8, 0)
    {Controladora\\de disco};
  \node[bloco es, minimum width=34mm, minimum height=13mm] (cr) at (13.2, 0)
    {Controladora\\de rede};

  % --- A faixa do barramento ---------------------------------------------
  \def\esq{-1.6}
  \def\dir{14.8}
  \fill[fill=black!4, rounded corners=4pt]
    (\esq, -3.9) rectangle (\dir, -1.6);

  \draw[barramento, draw=dadoborda,
        {Stealth[length=3mm]}-{Stealth[length=3mm]}]
    (\esq + 0.3, -2.0) -- (\dir - 0.3, -2.0);
  \draw[barramento, draw=procborda, -{Stealth[length=3mm]}]
    (\esq + 0.3, -2.75) -- (\dir - 0.3, -2.75);
  \draw[barramento, draw=esborda,
        {Stealth[length=3mm]}-{Stealth[length=3mm]}]
    (\esq + 0.3, -3.5) -- (\dir - 0.3, -3.5);

  \node[rotulo peq, anchor=west, text=dadoborda] at (\dir + 0.2, -2.0)
    {dados\\{\scriptsize o valor lido ou escrito}};
  \node[rotulo peq, anchor=west, text=procborda] at (\dir + 0.2, -2.75)
    {endereços\\{\scriptsize onde ler ou escrever}};
  \node[rotulo peq, anchor=west, text=esborda] at (\dir + 0.2, -3.5)
    {controle\\{\scriptsize ler, escrever, interromper}};

  \node[rotulo peq, anchor=north west] at (\esq, -4.05)
    {\textbf{Barramento do sistema}};

  % --- Ligação de cada bloco à faixa -------------------------------------
  \foreach \n in {cpu, mem, cd, cr} {
    \draw[seta dupla, line width=1.2pt] (\n.south) -- (\n.south |- 0,-1.6);
  }

  \node[rotulo peq, anchor=north east] at (\dir, -4.05)
    {O barramento de endereços tem sentido único, do processador para os
     demais.};

\end{tikzpicture}
\end{document}
